\subsection{Coupled-Mamba Fusion}

The coupled-Mamba fusion module is the core component used to integrate temporal, visual, and text-compatible representations in Time-ME. Given the temporal representation from the temporal branch, the visual embedding produced by the reconstruction-oriented MAE encoder, and the text-compatible embedding, we first project all modalities into a unified latent space through modality-specific projection layers. This alignment allows the subsequent state-space blocks to process heterogeneous inputs within a common feature dimension.

Let the aligned temporal, visual, and text features be denoted by $T_0$, $V_0$, and $S_0$, respectively. These features are then fed into three parallel modality-specific backbones, each consisting of a stack of Mamba-based state-space blocks. Unlike a standard early-fusion strategy, the three branches preserve modality-specific dynamics during the main sequence modeling stage. In our implementation, the temporal branch operates on variable-wise tokens, while the visual and text branches are short sequences and therefore introduce only modest computational overhead.

Formally, let $T \in \mathbb{R}^{B \times n \times d}$, $V \in \mathbb{R}^{B \times 1 \times d_v}$, and $S \in \mathbb{R}^{B \times d_s}$ denote the temporal, visual, and text inputs to the fusion module. After modality-specific projections, we obtain
\begin{gather}
T_0 = g_t(T), \qquad V_0 = g_v(V), \qquad S_0 = g_s(S),
\end{gather}
where $g_t(\cdot)$, $g_v(\cdot)$, and $g_s(\cdot)$ map the three modalities into the same latent space $\mathbb{R}^{d}$. For each layer $l \in \{1,\dots,L\}$, the modality-specific Mamba blocks follow a residual-normalized update:
\begin{gather}
\tilde{H}_m^{(l)} = \operatorname{Mamba}_m^{(l)}\!\left(H_m^{(l-1)}\right), \qquad
H_m^{(l)} = \operatorname{Norm}\!\left(\tilde{H}_m^{(l)}\right) + H_m^{(l-1)},
\end{gather}
where $m \in \{t,v,s\}$ denotes the temporal, visual, and text branches, and $H_t^{(0)} = T_0$, $H_v^{(0)} = V_0$, $H_s^{(0)} = S_0$.

To capture complementary information across modalities, we further introduce a cross-modal enhancement step after each Mamba layer. Small residual MLPs inject information from the other two branches into each modality. The temporal branch is enhanced by broadcasting the visual and text features to the temporal token length, while the visual and text branches are enhanced by the pooled temporal representation:
\begin{gather}
T^{(l)} = H_t^{(l)} + R_t\!\left([\operatorname{Broadcast}(H_v^{(l)});\operatorname{Broadcast}(H_s^{(l)})]\right), \\
V^{(l)} = H_v^{(l)} + R_v\!\left([\operatorname{Mean}(H_t^{(l)}); H_s^{(l)}]\right), \\
S^{(l)} = H_s^{(l)} + R_s\!\left([\operatorname{Mean}(H_t^{(l)}); H_v^{(l)}]\right),
\end{gather}
where $R_t(\cdot)$, $R_v(\cdot)$, and $R_s(\cdot)$ are lightweight residual enhancement networks, $[\cdot;\cdot]$ denotes concatenation, and $\operatorname{Mean}(\cdot)$ denotes sequence-wise average pooling of the temporal branch.

At the final layer, each branch output is pooled along its sequence dimension to obtain compact modality summaries $p_t$, $p_v$, and $p_s \in \mathbb{R}^{B \times d}$. A set of learnable scalar parameters is normalized by a softmax operator to produce modality weights, and the final fused representation is computed as
\begin{gather}
p_m = \operatorname{MaxPool}\!\left(H_m^{(L)}\right), \qquad
\alpha = \operatorname{Softmax}([w_t,w_v,w_s]), \\
f = \sum_{m \in \{t,v,s\}} \alpha_m p_m, \qquad
Y_{mm} = h_{\mathrm{reg}}(f),
\end{gather}
where $h_{\mathrm{reg}}(\cdot)$ denotes the regression head. The multimodal prediction $Y_{mm}$ is then combined with the memory-branch prediction through a learnable sigmoid gate in the final forecasting stage. In this way, the coupled-Mamba module provides a multimodal forecast that preserves modality-specific dynamics while explicitly modeling cross-modal interactions.
